\documentclass[
    a4paper,
    pagesize,
	pdftex,
    12pt,
]{scrartcl}

\usepackage[T1]{fontenc}
\usepackage[ngerman]{babel}

\usepackage[unicode=true]{hyperref}
\usepackage{graphicx}
\usepackage[nopatch=footnote,draft=false,babel,tracking=true,kerning=true,spacing=true]{microtype}

\usepackage{enumerate}
\usepackage{subcaption}
\usepackage{abstract}
\usepackage{amsmath}
\usepackage{mathtools}
\usepackage{listings}
\usepackage{enumitem}
\usepackage{longtable}
\usepackage[table]{xcolor}  

\graphicspath{{./pictures/}}

\begin{document}

\begin{titlepage}
    \begin{center}
        
        \includegraphics[height=25mm]{S04_HTW_Berlin_Logo_pos_FARBIG_RGB.jpg} \\
        \vspace{1.0cm}
         Multi-Level Erkennung von IoT-Botnetzen in frühen Botnetzphasen
             
        \vspace{1.5cm}
 
        \textbf{Exposé zur Dissertation} \\    
        zur Erlangung des Grades eines\\
        Doktors

        \vspace{1.5cm}
        am Fachbereich 4 - Informatik, Kommunikation und Wirtschaft\\
        der Hochschule für Technik und Wirtschaft Berlin\\
        
        \vspace{1.5cm}

        vorgelegt von \\
        \textbf{Jan Schröder}
      
        %\includegraphics[width=0.4\textwidth]{university}
        \vspace{1.5cm}    
        Berlin, \today\\
        
             
    \end{center}
\end{titlepage}

\pagenumbering{gobble}

\tableofcontents
\newpage

\pagenumbering{arabic}

% left fixed width:
\newcolumntype{L}[1]{>{\raggedright\arraybackslash}p{#1}}

\begin{longtable}[c]{ |L{4cm}|L{10cm}| }
    \hline
    \rowcolor{lightgray} Begriff & Erklärung \\ [0.5ex]
    \hline
    \endfirsthead

    Internet of Things & Geräte, zum Beispiel aus dem Haushalt, die mit dem Internet verbunden werden können \cite{NIST800171}. \\
    \hline
    Malware & Software, die durch Ausnutzen einer Schwachstelle zu einem schädlichen Zweck führt \cite{CNSSI4009}. \\
    \hline
    Honeypot & Computersystem um Informationen eines Angreifers zu bekommen, indem es gezielt kompromittiert wird \cite{DBLP:conf/icccs/MairhBVJ11}. \\
    \hline
    Software Defined Network & Architektur, die die Netzwerksteuerung von den Routern und Switches trennt, welche den Netzwerkverkehr weiterleiten \cite{DBLP:journals/pieee/KreutzRVRAU15}. \\
    \hline
    Fog Computing & Bonomi et al. \cite{DBLP:conf/sigcomm/BonomiMZA12} stellen die Fog Computing Architektur vor. \\
    \hline
    Edge Computing & Edge Computing ist eine Architektur die zum Fog Computing zählt \cite{DBLP:journals/fgcs/KhanAHYA19}. \\
    \hline
    Network Function Virutalization & Konzept, um verschiedene Funktionen virtualisiert auf einem Server auszuführen. So wird nicht mehr ein Server für eine Funktion benötigt, wie zum Beispiel eine Firewall, die auf einem Server ausgeführt wird \cite{DBLP:journals/corr/MijumbiSGBTB15}. \\
    \hline
    \caption{Diese Tabelle erläutert Begriffe, welche sowohl im Exposé als auch der späteren Dissertation Verwendung finden sollen.\label{tab:term_table}}\\
\end{longtable}
\newpage
\section{Forschungsfragen und Aufbau des Exposés}

Botnetze im Internet of Things (IoT) Gebiet haben in der modernen Welt sehr stark an Bedeutung gewonnen. Gerade im privaten 
Haushalt finden sich immer mehr IoT-Geräte die durch Botnetze wie Mirai und Hajime infiziert werden können. Daher ist es 
sehr wichtig, dass diese Botnetze durch eine Erkennung identifiziert und abgewehrt werden. Die Dissertation soll sich mit 
der Erkennung von Botnetzen beschäftigen und die Identifikation optimieren.

\subsection*{Forschungsfragen für die Dissertation}

\newlist{questions}{enumerate}{2}
\setlist[questions,1]{label=\textsf{\textbf{Q\arabic*:}},ref=RQ\arabic*}
\setlist[questions,2]{label=(\alph*),ref=\thequestionsi(\alph*)}

Die Dissertation soll spezifisch die folgenden Forschungsfragen beantworten:  

\begin{questions}
    \item Wie lassen sich Botnetze während der Verbreitungsphase erkennen? \label{itm:q1}
    \item Wie lassen sich Botnetze während der Scan-Phase erkennen? \label{itm:q2}
    \item Welche Kombinationen von bereits bekannten Methoden führen zu einer Verbesserung der Erkennung von Botnetzen? \label{itm:q3}
    \item Welche Kombinationen von bereits bekannten Techniken führen zu einer Verbesserung der Erkennung von Botnetzen? \label{itm:q4}
    \item Wie lässt sich die Erkennung von Botnetzen optimieren, wenn mehrere Erkennungsstufen verwendet werden? \label{itm:q5}
\end{questions}

In Frage \ref{itm:q1} und \ref{itm:q2} soll es um die Erkennung von Botnetzen während der ersten beiden Phasen gehen. Was die Phasen sind erklärt das nachfolgende Kapitel. 
\\ \\
In den folgenden Kapitel \ref{sec:theory} werden die theoretischen Hintergründe zum Thema der Dissertation und verschiedene Prozesse zur Erkennung von Botnetzen, sowie 
mehrere Botnetze die in Fallstudien eingesetzt werden sollen beschrieben. Kapitel \ref{sec:research} erläutert den aktuellen Stand der Forschung zur Erkennung von Botnetzen,
sowie zu den Themen die in der Dissertation behandelt werden sollen. In Kapitel \ref{sec:goals} werden die Ziele der Dissertation aufgestellt. Das Kapitel \ref{sec:methods}
erklärt wie in der Dissertation vorgegangen werden soll, um ein Konzept mit einer dazugehörigen Implementierung zu erstellen sowie den geplanten Aufbau eines Laborexperiments
und mögliche Fallstudien. Das letzte Kapitel \ref{sec:plan} stellt einen Zeitplan vor sowie die Struktur, die die Dissertation haben soll.
\newpage
\section{Einführung in die Arbeit}

\todo[inline]{Wie unterscheidet sich die Arbeit von anderen (mit Beispielen). Warum ist das Thema notwendig. Es muss Drama eingebracht werden, warum die Arbeit so wichtig ist. Es muss den Leser motivieren.}

\subsection{Motivation}

\subsection{Eingrenzungen}



\newpage
\section{Stand der Forschung}
\label{sec:research}

\subsection*{Methoden zur Erkennung von Botnetzen}
% - SDN, Edge Computing etc. Referenz muss hierhin und einzelne Techniken noch erklärt werden
Xing et al. \cite{Xing2021SurveyOB} stellt eine zusammengefasste Unterteilung von Methoden zur Botnetzerkennung vor. Dabei gehen sie aber nicht konkreter auf Methoden
ein. Speziell auf Botnetze im IoT Gebiet fassen Wazzan et al. \cite{Wazzan2021InternetOT} Methoden und Techniken anhand von Studien zusammen, die zur Erkennung 
eingesetzt werden. Für den aktuellen Stand der Forschung nutzt dieses Exposé deshalb die ausgearbeiteten Zusammenfassungen von \cite{Wazzan2021InternetOT} um einen 
Ansatz zu finden, wie die Botnetzerkennung stattfinden könnte.
Im Zusammenhang mit der Botnetzerkennung findet ein Großteil der Forschung zu diesem Gebiet im Zusammenhang mit Maschinellem Lernen (ML) statt, wie unter anderem zum Beispiel
\cite{SAlrayes2022ModelingOB,DBLP:journals/comcom/Alani22,Habtamu2022ASR, D2022ARO} zeigen. Wazzan et al. schlägt die Kombination 
mehrerer Erkennungsmethoden vor, womit sich die Dissertation auch beschäftigen soll. Nach verschiedener Studien schlagen Wazzan et al. folgende Erkennungsmethoden vor: 
Software Defined Network (SDN), Edge Computing, Blockchain, Fog Computing und Network Function Virtualization (NFV). Im Folgenden wird der aktuelle Stand der Forschung zu jeder dieser 
Technologien betrachtet. \\ \\ Zha et al. \cite{DBLP:conf/cns/Zha0GMC19} beschreiben eine Bibliothek, welche Botnetz Erkennung anhand eines SDNs durchführt. 
Bei dem Konzept des Frameworks handelt es sich darum, Botnetz und C\&C-Kanal Aktivitäten zu identifizieren, damit die Robustheit der Erkennung erhöht wird. Zum Aufbau gehört ein SDN 
Controller, welcher ein ML Modell enthält, dass über Erkennung und Überwachung trainiert. Die Überwachung des Netzwerks, findet über Software Switches statt, welche vollen Zugang zum 
Netzwerkverkehr von Virtuellen Maschinen (VM) haben. Das ML Modell wird für das komplette Netzwerk verwendet, um Bot Aktivitäten zu erkennen, während für einzelne Server, dessen VMs mit 
einem Software Switch kommunizieren, lokale Netzwerküberwachung stattfindet. Über die einzelnen Software Switches, findet die C\&C Identifizierung statt. Dabei ist speziell Peer-to-Peer 
Verkehr, HTTP und IRC Protokolle im Fokus. Das ML Modell ist ein Neuronales Netzwerk (NN), welches kompromittierte Hosts findet anhand von aufgebauten Verbindungen des Hosts. Die Entscheidung, 
ob ein Host zu einem Bot wird, entscheidet das NN über einen festgelegten Schwellenwert. Mit einer einzelnen positiven Identifizierung ist aber dennoch nicht eindeutig geklärt, dass der
Host auch wirklich ein Bot ist, da er zum Beispiel nicht unbedingt mit dem C\&C-Server kommuniziert. An dieser Stelle wird der Host dennoch dauerhaft weiter beobachtet. Der aktuelle Stand Forschung 
zeigt hier, dass die Erkennung nicht während der Scan-Phase durchführt wird, womit sich ein Ziel ergibt Botnetzerkennung in frühen Botnetzphase durchzuführen, was in der Dissertation verfolgt 
werden soll und weiter in Kapitel \ref{sec:goals} erläutert wird. Zusätzlich zu dem beschriebenen Framework erläutern Negera et al. \cite{DBLP:journals/sensors/NegeraSDMA22} verschiedene Arbeiten 
zur Erkennung von Botnetzen in SDNs über Modelle des ML. \\ \\ Neben der Erkennung mit SDNs findet in der aktuellen Forschung auch die Erkennung von Botnetzen über Edge Computing Architekturen 
statt, wie \cite{gromov2022edge} in ihrer Arbeit zeigen. Bei dieser Erkennung verwenden Gromov et al. in einem Edge Computing System ein Convolutional Neural Network (CNN). Der Netzwerkverkehr der 
einzelnen IoT-Geräte wird an ein NVIDIA Jetson Nano Entwickler Kit zur Botnetzerkennung geleitet. Dabei werden anhand von Computer Vision Bilder zur Erkennung erzeugt, welche in auffällige und 
unauffällige Bilder klassifiziert werden, um die Performance zu erhöhen sowie das Netzwerk zu verkleinern. \\ \\ Bei der Blockchainbasierten Erkennung geht es darum, eine Blockchain zu implementieren 
und in diesem Netzwerk ein IoT Netzwerk auf Botnetze zu überwachen. Dabei werden unterschiedliche Methoden innerhalb der Blockchain implementiert und mit den Vorteilen der Blockchain verwendet. 
Sagirlar et al. \cite{DBLP:journals/corr/abs-1809-10775} nutzen diese Methode indem sie innerhalb der Blockchain eine Peer-to-Peer Botnetz Erkennung implementieren. Salim et al. 
\cite{DBLP:journals/sensors/SalimATPP22} erläutern ein Framework, welches Digital Twins und Packet Auditor in einer Blockchain zur Erkennung einsetzen. \\ \\ Neben Blockchains wird auch Fog 
Computing zur Botnetzerkennung eingesetzt. \cite{Lawal2020AnAM} stellen ein Framework vor, welches Anomalien in einem IoT Netzwerk erkennt, aufgebaut als Fog Computing Architektur, anhand von signatur- 
und anomaliebasierten Erkennungsmethoden. Die signaturbasierte Methode greift auf eine Datenbank zu, welche IP-Adressen enthält die als Angriffsursprung identifiziert sind. Die anomaliebasierte Methode 
klassifiziert anhand eines Extreme Gradient Booster (XGBoost) Algorithmus \cite{DBLP:journals/corr/ChenG16} zwischen auffälligem und unauffälligem Netzwerkverkehr. Der Netzwerkverkehr fließt zu Beginn 
durch ein signaturbasiertes IDS, um die IP-Adresse mit der Datenbank zu vergleichen. Sollte nichts gefunden werden, fließt der Verkehr durch ein anomaliebasiertes IDS und klassifiziert ihn entsprechend.
\\ \\ Das letzte Framework erläutern Kim et al. \cite{DBLP:journals/cn/KimNPSS19} in ihrer Arbeit zu einer SDN Umgebung mit integrierten Netzwerk-Funktionen auf NFV, welche in den IoT-Geräten integriert
sind. Das Konzept aus der Arbeit von Kim et al. beschreibt ein IoT Netzwerk, welches für jede einzelne Komponente eigene Sicherheitsrichtlinien vorgibt. Die Architektur des Netzwerks beginnt mit einer
\textit{Control Plane}, welches für das Management des Netzwerks verantwortlich ist. Die \textit{Control Plane} besteht aus fünf Komponenten: \textit{Datapath}, welches für die Kommunikation zwischen 
\textit{Control Plane}, \textit{Function Plane} und \textit{Data Plane} (Verbindung zu den IoT-Geräten) verantwortlich ist. Die zweite Komponente ist der \textit{Protocol Parser} der Nachrichten von der 
\textit{Function-} und \textit{Data Plane} empfängt und zum \textit{Core Module} weiterleitet. Das \textit{Core Module} verwaltet Netzwerkkomponenten, stellt 
Sicherheitsfunktionen zur Verfügung und erkennt Anomalien bei den Sicherheitsrichtlinien. Der \textit{Event Manager} stellt Event basierte Kommunikation zwischen zur Verfügung. Die letzte Komponente sind 
\textit{Applications} welche über ein \textit{programmable interface} implementiert sind und sind für die Verarbeitung von den IoT-Geräten zuständig. Die \textit{Function Plane} ist für die Sicherheit der 
IoT-Geräte zuständig, da diese durch begrenzte Leistung anspruchsvolle Sicherheitsfunktionen nicht verarbeiten können. Diese Sicherheitsfunktionen werden über NFV Techniken ausgeführt, damit die Funktionen 
virtualisiert zur Verfügung stehen. Die Kommunikation findet über das von Kim et al. vorgestellte SODA Protocol statt, damit zum Beispiel Richtlinien richtig verwaltet werden können. Wenn zum Beispiel neue 
Geräte dem Netzwerk hinzugefügt werden, dann werden anhand von eventbasierten Prozessen diese auf zum Beispiel Sicherheitsbedenken geprüft. Sollte über das Netzwerk ein Angriff stattfinden, so wird das an 
den Richtlinien erkannt und diese dann angepasst. \\ Bei der Implementierung wurde die \textit{Control-} und \textit{Function Plane} als SDN Controller implementiert, das \textit{Data Plane} als Kommunikation 
zwischen dem Controller und den IoT-Geräten. Zur vereinfachten Ressourcenverwaltung nutzen Kim et al. NFV bei jedem IoT-Gerät. \\ \\ Aus den Arbeiten der aktuellen Forschung geht hervor, dass bei der Verwendung 
von ML die eingesetzten Algorithmen immer zu Ergebnissen führen bei denen das Modell eine Genauigkeit von über 90\% erreicht. Dabei gehen die Arbeiten nicht weiter darauf ein, warum der entsprechende Algorithmus
verwendet wird und welche Vorteile dieser bringt. Zudem erkennen die Arbeiten die Botnetze immer zu den späteren Phasen eines Botnetzes womit die Arbeiten keine Prävention von Botnetzen durchführen. Des weiteren
nutzen die Botnetzerkennungen nur einzelne Methoden und kombinieren diese nicht miteinander um zum Beispiel die Nachteile einer Methode mit einer weiteren auszugleichen. Wenn die Netzwerkarchitekturen wie 
zum Beispiel SDNs mit in die Erkennung von IoT Botnetzen mit einbezogen werden, zeigt die aktuelle Forschung, dass dies zu einer besseren Leistung der Botnetzerkennung führt, um so zum Beispiel Latenzen zu verringern.
Aus Kapitel \ref{sec:theory} und dem aktuellen Stand der Forschung lassen sich \cite{Xing2021SurveyOB} und \cite{Wazzan2021InternetOT} als mögliche Ansätze für die Botnetzerkennung verwenden.

\newpage
\section{Ziel der Dissertation}
\label{sec:goals}

Mit der Dissertation soll die Botnetzerkennung optimiert werden. Ziel ist es, ein Konzept zu erstellen, zusammen mit einer Implementierung.
Das erste Ziel soll ein Konzept zur Botnetzerkennung werden, welches zu jeder Phase eines Botnetzes, mehrere Erkennungslevel integriert, die aus einer Kombinationen 
mehrerer Erkennungsmethoden besteht. Der aktuelle Stand aus der Forschung zu den einzelnen Erkennungsmethoden zeigt außerdem, dass die aktuellen Methoden sich eher 
auf die späteren Phasen eines Botnetzes fokussieren. Daher ist ein weiteres Ziel der Fokus auf die Entwicklung zur Erkennung der früheren Phasen der Botnetze.
Das Ziel der Erkennungsmethoden ist es die Methoden herauszuarbeiten, welche die wenigsten false positives erzielen um diese dann miteinander zu kombinieren. Welche 
Methoden genau eingesetz werden, soll mit dieser Dissertation zu Beginn herausgefunden werden. Ein weiteres Ziel ist, ob die Erkennung auch unbekannte Botnetze 
identifizieren kann indem Fallstudien durchgeführt werden. \\ Ziel der Dissertation ist es nicht, aus der Erkennung Maßnahmen zur Verteidigung gegen Botnetze zu treffen. 
Die Dissertation soll sich ausschließlich auf die Erkennung konzentrieren. 
\newpage
\section{Forschungsdesign und Methodik}
Die Dissertation geht mit einem Laborexperiment einher. Um die Erkennung eines 
Botnetzes zu testen sollen verschieden typische IoT-Geräte aufgebaut werden, die 
von Botnetzen wie Mirai angegriffen werden. Dazu ist es nötig, dass Geräte mit einem
ARC-Prozessor aufgebaut sind. \\ \\
In einer Fallstudie ist dann vorgesehen, das Botnetz auszuführen,
Daten den Erkennungsindikatoren zuzuweisen und dann durch ein ML Modell das Botnetz erkennen.

\newpage
\section{Zeitplan und Struktur der Dissertation}
\label{sec:plan}

Die voraussichtliche Dauer der Dissertation beträgt 4 Jahre und 6 Monate.

\begin{center}
    \begin{tabular}[c]{| c | c |}
        \hline
        \rowcolor{lightgray} Zeit & Vorgehen \\ [0.5ex]
        \hline
        März 2023 & Grundlagen, verwandte Arbeiten erläutern \\
        \hline
        August 2023 &  Laborexperiment starten, Konzept in der Dissertation beschreiben \\
        \hline
        Januar 2025 & Implementierung der IoT Geräte beschreiben \\
        \hline
        März 2025 & Fallstudien beschreiben \\
        \hline
        Juli 2025 & Ergebnisse erläutern, Fazit schreiben \\
        \hline
        Februar 2026 & Einleitung schreiben, Kontrolle der Dissertation \\
        \hline
        Juli 2026 & Abgabe \\
        \hline
    \end{tabular}
\end{center}

Zum Ende der Arbeitszeit soll die Dissertation über eine externe Stelle auf Plagiate, Rechtschreibung etc.
geprüft werden.

\subsection*{Kapitelstruktur der Dissertation}

\begin{enumerate}[label=\Roman*]
    \item Einführung in die Dissertation
    \item Grundlagen und Verwandte Arbeiten
    \item Konzept der Mehrstufigen Botnet Erkennung
    \item Implementierung des Experiments und der Fallstudien
    \item Evaluierung und Ergebnisse
\end{enumerate}

% - In der Struktur sollen die Forschungsfragen als eigene Kapitel beantwortet werden und die Ergebnisse 
%   strukturiert als Erstens, zweitens,...etc.!!!!
\newpage

\bibliographystyle{ieeetr}
\bibliography{references}

\end{document}
